% !TeX program = xelatex
% UTF-8 transcription of all 17 original photographs.
% Compile this file directly with XeLaTeX, or input it into a XePersian book.
% Original spelling, abbreviations and medical statements are retained.
% Source-photo numbering follows ascending original filenames.
\newif\ifMotorStandalone
\ifdefined\settextfont\MotorStandalonefalse\else\MotorStandalonetrue\fi
\ifMotorStandalone
  \documentclass[a4paper,12pt]{extbook}
  \usepackage[margin=2.2cm]{geometry}
  \usepackage{xcolor}
  \usepackage{amsmath,amssymb}
  \usepackage{enumitem}
  \usepackage{xepersian}
  \settextfont{Amiri}
  \setlatintextfont{FreeSerif}
\fi
\ifMotorStandalone\begin{document}\fi
\begingroup
\emergencystretch=3em
\newcommand{\MDen}[1]{\lr{#1}}
\newcommand{\MDunclear}[1]{[خوانش نامطمئن: #1]}
\newcommand{\MDhl}[1]{\ifdefined\colorbox\colorbox{yellow!35}{\strut #1}\else\textbf{#1}\fi}
\newcommand{\MDpage}[3]{\clearpage\section*{صفحهٔ #1}\noindent\MDen{\texttt{#2}}\par\noindent موضوع: #3\par\medskip
\noindent{\footnotesize\MDen{Date: / / \quad Subject: \quad Sat. Sun. Mon. Tue. Thu. Wed. Fri.}}\par\medskip}
\chapter{اختلالات حرکتی}

% pic01.jpg = 20261012_224433637.jpg; handwritten page 1.
\MDpage{۱}{pic01.jpg}{۱۴ ـ اختلالات حرکتی}
\begin{itemize}
\item بازال گانگلیا سبب اختلالات حرکتی می‌دهد.
\item پوتامن، هستهٔ کودیت، گلوبوس پالیدوس، سابستانتیا نیگرا.
\begin{itemize}
\item پوتامن و هستهٔ کودیت: جسم استریاتوم.
\item گلوبوس پالیدوس: پالئو استریاتوم.
\end{itemize}
\end{itemize}
\section*{پارکینسون}
\begin{itemize}
\item سن: \MDunclear{دههٔ ۴–۵۰ سال} / شیوع با بالا رفتن سن \MDen{$\uparrow$} / \MDen{$M=F$}.
\item در غیرسیگاری‌ها شایع‌تر.
\item پاتوژنز: استرس اکسیداتیو، تخریب نورون‌های پیگمانتهٔ سابستانتیا نیگرا.
\begin{itemize}
\item مواد توکسیک مثل \MDen{MPTP} (متیل فنیل تتراهیدروپیریدین): علائم سریعاً پیشروندهٔ شبه‌پارکینسون.
\item ژنتیک.
\end{itemize}
\item پاتولوژی: از بین رفتن نورون‌های دوپامینرژیک پیگمانته در جسم سیاه.
\begin{itemize}
\item وجود اجسام لوی (انکلوزیون‌های ائوزینوفیلیک): حاوی پروتئین \MDen{$\alpha$-Synuclein}.
\end{itemize}
\item کرایتریای تشخیص: وجود ۳ معیار از کرایتریاهای تشخیص:
\begin{enumerate}[label=\alph*)]
\item وجود ۲ علامت از ۳ علامت: برادی‌کینزی / ترمور در حالت استراحت / ریجیدیتی و سفتی عضلانی.
\item پاسخ بالینی به لوودوپا.
\begin{itemize}
\item بیماری پارکینسون در تمام مراحل به لوودوپا پاسخ می‌دهد.
\item پاسخ به دوز \MDen{1.5 g} روزانه، به مدت (حداقل ۳ ماه).
\end{itemize}
\item عدم وجود علائم زیر:
\begin{itemize}
\item علائم مخچه‌ای / علائم پیرامیدال / افتالموپلژی و محدودیت حرکات چشم.
\item سقوط مکرر در اوایل بیماری / آپراکسی / دیسفاژی و دیزآرتری زودرس.
\item دیسفاژی و استریدور / علائم زودرس و شدید اتونوم / دمانس زودرس.
\end{itemize}
\end{enumerate}
\item توهمات بینایی با دوز اندک داروهای ضدپارکینسون.
\end{itemize}

% pic02.jpg = 20261012_224440277.jpg; handwritten page 2.
\MDpage{۲}{pic02.jpg}{۱۴ ـ اختلالات حرکتی}
\section*{تظاهرات نورولوژیک پارکینسون}
\subsection*{تظاهرات آکینتیک}
\begin{enumerate}[label=\alph*)]
\item برادی‌کینزی: کندی حرکات.
\item هایپوکینزی: کاهش دامنهٔ حرکات.
\item آکینزی: اشکال در شروع حرکت.
\end{enumerate}
\begin{itemize}
\item مهم‌ترین و ناتوان‌کننده‌ترین علامت پارکینسون.
\end{itemize}
\subsection*{ترمور در حین استراحت}
\begin{itemize}
\item در \MDen{$\tfrac{2}{3}$} بیماران اولین علامت است.
\item در نهایت همهٔ بیماران دچار آن می‌شوند.
\item ترمور حین استراحت و با \MDen{$f=3$--$6\,\mathrm{Hz}$}.
\item ترمور در دست شایع و به صورت \MDen{Pill-Rolling} (تسبیح چرخاندن).
\item با استرس روانی تشدید / در خواب عمیق رفع می‌شود.
\item در پا و فک هم ممکن است رخ دهد.
\end{itemize}
\subsection*{ریجیدیتی عضلانی}
\begin{itemize}
\item مقاومت در تمام طول \MDen{ROM} عضلات آگونیست و آنتاگونیست.
\item ۲ نوع ریجیدیتی شایع در \MDen{PD}:
\begin{enumerate}[label=\alph*)]
\item لوله‌سربی (\MDen{Lead-pipe rigidity}).
\item چرخ‌دنده‌ای (\MDen{Cogwheel-rigidity}): ریجیدیتی همراه با ترمور.
\end{enumerate}
\end{itemize}
\subsection*{وضعیت خمیده و ناپایداری وضعیت}
\begin{itemize}
\item راه رفتن \MDen{PD} بسیار اختصاصی است.
\item آهسته در وضعیت خمیده به جلو با قدم‌های کوتاه و کشیدن پاها روی زمین.
\end{itemize}

% pic03.jpg = 20261012_224452656.jpg; handwritten page 3.
\MDpage{۳}{pic03.jpg}{۱۴ ـ اختلالات حرکتی}
\begin{itemize}
\item حرکات طبیعی دست حین راه رفتن \MDen{$\downarrow$}: کاهش \MDen{Arm Swing}.
\item فلکشن خفیف در مفصل آرنج، زانو، ران.
\item \MDen{Festinating gate}: پدیدهٔ تندروی، بسیار کلاسیک و اختصاصی و کاراکتریستیک.
\begin{itemize}
\item بیمار تمایل به سقوط به جلو دارد، لذا هنگام راه رفتن با هر قدمی که برمی‌دارد برای جلوگیری از سقوط قدم بعدی را سریع‌تر برمی‌دارد. راه رفتن با قدم‌های کوتاه و سریع و با قامت خمیده.
\end{itemize}
\item در مراحل پیشرفته: توقف در شروع راه رفتن یا چرخیدن.
\begin{itemize}
\item انجماد حرکتی: \MDen{Freezing}: بلوک موتور.
\end{itemize}
\end{itemize}
\subsection*{سایر علائم نورولوژیک}
\begin{itemize}
\item صورت بی‌حالت و ماسکه.
\item میکروگرافی: دست‌خط بد و ریز می‌شود.
\item صدا آهسته، کوتاه و مونوتون می‌شود.
\item بابنسکی \MDen{Downward}، \MDen{DTR} نرمال است.
\item دمانس و اختلالات ذهنی در مراحل پیشرفته.
\item برادی‌فرنی: کندی روند تفکر و اختلال در یادآوری.
\item معاینهٔ حسی طبیعی است ولی علائم حسی در آن‌ها شایع می‌باشد:
\begin{itemize}
\item درد اندام‌ها / بی‌قراری حسی پاها (آکاتیزیا) / احساس درد و سوزش اندام‌ها.
\end{itemize}
\end{itemize}
\subsection*{علائم غیرنورولوژیک}
\begin{itemize}
\item خستگی: شایع / ژنرالیزه یا محدود به یک اندام.
\item گوارشی: یبوست / کاهش وزن.
\item در مراحل پیشرفته: دیسفاژی و ریزش بزاق از دهان.
\item چشمی: بلفارواسپاسم، بلارویژن، محدودیت نگاه به بالا، فلج \MDen{Convergence}.
\begin{itemize}
\item حرکات ساکادیک و تعقیبی چشم \MDen{N/L}.
\end{itemize}
\end{itemize}

% pic04.jpg = 20261012_224506162.jpg; handwritten page 4.
\MDpage{۴}{pic04.jpg}{۱۴ ـ اختلالات حرکتی}
\begin{itemize}
\item اختلالات ادراری: تکرر، فوریت، بی‌اختیاری \MDen{Urge}: تحریک‌پذیری مثانه.
\item هایپوتنشن وضعیتی: دیس‌اتونومی / مصرف داروهای دوپامینرژیک.
\item افسردگی: در \MDen{$\tfrac{1}{3}$} بیماران / شدید نیست و از نوع واکنشی می‌باشد.
\end{itemize}
\section*{مطالعات تشخیصی در \MDen{PD}}
\begin{itemize}
\item \MDen{MRI} مغز \MDen{N/L} می‌باشد.
\item \MDen{SPECT}: اختلال در سیستم دوپامینرژیک گانگلیون بازال؛ افتراق \MDen{PD} از ترمور اسنشیال.
\item ترمور، ریجیدیتی و برادی‌کینزی در \MDen{PD} به صورت غیرقرینه است.
\item \MDen{PD} معمولاً از اندام فوقانی یک طرف آغاز می‌شود؛ به‌ندرت از اندام تحتانی آغاز می‌شود.
\item حتی در فرم شدید \MDen{PD}، غیرقرینگی باقی می‌ماند.
\item اگر علائم دمانس بعد از یک سال از شروع \MDen{PD} آغاز شود به آن پارکینسون به همراه دمانس (\MDen{PDD}) می‌گویند.
\item اگر در کمتر از یک سال آغاز شود به آن دمانس با لوی بادی گفته می‌شود.
\item در هر دو توهم بینایی شایع است ولی اختلال توجه، تمرکز، اختلال سطح هوشیاری و عملکرد شناختی همراه با نوسان از ویژگی دمانس با لوی بادی است.
\item در دمانس با لوی بادی، اختلال شناخت، عفونت ادراری یا تنفسی منجر به دلیریوم می‌شود.
\end{itemize}

% pic05.jpg = 20261012_224514045.jpg; handwritten page 5.
\MDpage{۵}{pic05.jpg}{۱۴ ـ اختلالات حرکتی}
\section*{درمان پارکینسون}
\begin{itemize}
\item اندیکاسیون درمان: در مراحل ابتدایی بیماری و با علائم خفیف درمان دارویی لازم نیست.
\item زمانی که علائم زندگی روزمره را مختل کنند درمان می‌کنیم.
\item انتخاب دارو: بر اساس سن و وضعیت فعال.
\item بیماران جوان: شروع با آگونیست دوپامین مثل پرگولید و بروموکریپتین؛ به صورت نیاز لوودوپا اضافه می‌شود.
\item بیمار مسن: درمان با لوودوپا شروع می‌شود و بهتر است داروهای دیگر اضافه نشود.
\end{itemize}
\subsection*{لوودوپا}
\begin{itemize}
\item مؤثرترین دارو برای تمام علائم پارکینسون.
\item مؤثرترین دارو در کنترل ترمور.
\item باید به صورت ترکیبی با کربی‌دوپا یا بن‌سرازید تجویز شود؛ برای اینکه خارج از \MDen{CNS} به دوپامین تبدیل نشود.
\item شروع با دوز \MDen{50--100 mg} روزانه بعد از غذا، سپس دوز دارو را \MDen{$\uparrow$} و مصرف دارو را به قبل از غذا منتقل می‌کنیم.
\item \MDen{PP} جذب لوودوپا را مختل می‌کند.
\item در اوایل بیماری اثر دارو وابسته به دوز نیست ولی در ادامه وابسته به دوز شده و به صورت تأخیر در مصرف، علائم پارکینسون ایجاد می‌شود.
\begin{itemize}
\item \MDen{End of dose deterioration}، \MDen{Wearing off}.
\item در نتیجه بیمار پدیدهٔ \MDen{ON/OFF} را تجربه می‌کند.
\end{itemize}
\end{itemize}
\subsection*{عوارض لوودوپا}
\begin{enumerate}[label=\alph*)]
\item دیس‌کینزی ناشی از لوودوپا.
\item عوارض اتونوم: شایع‌ترین: هایپوتنشن وضعیتی.
\item عوارض روانی: دلیریوم، توهم بینایی، سایکوز، دمانس.
\end{enumerate}

% pic06.jpg = 20261012_224522396.jpg; handwritten page 6.
\MDpage{۶}{pic06.jpg}{۱۴ ـ اختلالات حرکتی}
\noindent شمارهٔ حاشیهٔ چپ: ۴۳.
\subsection*{کنتراندیکاسیون‌های نسبی لوودوپا}
\begin{enumerate}[label=\alph*)]
\item نباید همراه با \MDen{MAOI-A} مصرف شود / حداقل ۱۵ روز از مصرف آن‌ها گذشته باشد.
\item \MDen{MI} اخیر.
\end{enumerate}
\subsection*{آگونیست‌های دوپامین}
\begin{itemize}
\item اندیکاسیون: مراحل اولیهٔ بیماری / افراد جوان / کنترل مشکلات ناشی از درمان با لوودوپا مثل دیس‌کینزی و پدیدهٔ \MDen{ON/OFF}.
\item پرگولید، بروموکریپتین، کابرگولین، لیزوراید، پرامی‌پکسول، روپینیرول، آپومورفین.
\item مهم‌ترین عارضه: \MDen{MI}، عوارض قلبی.
\end{itemize}
\subsection*{آمانتادین}
\begin{itemize}
\item مهم‌ترین کاربرد: درمان دیس‌کینزی ناشی از لوودوپا و دوپامین.
\end{itemize}
\subsection*{\MDen{MAOI-B}}
\begin{itemize}
\item سلژیلین، راساژیلین / این‌ها مؤثرتر است.
\item اثر اندک بر علائم پارکینسون.
\item مدت اثر لوودوپا را \MDen{$\uparrow$}.
\end{itemize}
\subsection*{آنتی‌کولینرژیک‌ها}
\begin{itemize}
\item تری‌هگزی‌فنیدیل، بی‌پریدن.
\item علیرغم اثر ضعیف بر علائم پارکینسون.
\item صرفاً روی ترمور و ریجیدیتی تا حدی اثر دارند.
\item مصرف در افراد مسن ممنوع است.
\item قطع ناگهانی آن‌ها موجب افزایش شدید علائم پارکینسون می‌شود.
\end{itemize}
\subsection*{جراحی}
\begin{itemize}
\item تالاموتومی، پالیدوتومی، تحریک عمقی مغز (\MDen{DBS}).
\end{itemize}

% pic07.jpg = 20261012_224529432.jpg; handwritten page 7.
\MDpage{۷}{pic07.jpg}{۱۴ ـ اختلالات حرکتی}
\section*{پیش‌آگهی پارکینسون}
\begin{itemize}
\item گروهی که در آن‌ها ترمور علامت غالب است پروگنوز بهتری دارند.
\item گروهی که با علائم آکینتیک، آکینزی و ریجیدیتی تظاهر می‌یابند پروگنوز بدتر.
\end{itemize}
\section*{فلج فوق‌هسته‌ای پیشرونده (\MDen{PSP})}
\begin{itemize}
\item یک اختلال دژنراتیو پیشرونده و نادر / \MDen{$F<M$} / سن ۴۵–۷۵ (متوسط ۶۳).
\end{itemize}
\subsection*{علائم بالینی}
\begin{itemize}
\item علامت اصلی افتالموپلژی فوق‌هسته‌ای است (اختلال در حرکات ارادی عمودی چشم به‌ویژه نگاه به پایین)؛ سپس فلج \MDen{gaze} در سایر جهات.
\item ریجیدیتی دیستونیک عضلات بالاتنه و گردن (قامت کشیده و گردن راست دارند) (به عقب کشیده شدن سر).
\item فلج سودوبولبار: اختلال پیشروندهٔ بلع و تکلم.
\item دمانس لوب فرونتال.
\item حملات سقوط مکرر به عقب (\MDen{Retropulsion}): کاراکتریستیک.
\item خنده و گریهٔ بی‌مورد.
\item رفلکس‌های اکولو سفالیک و اکولو وستیبولار \MDen{N/L} هستند.
\item بلفارواسپاسم.
\item علائم درگیری پیرامیدال: \MDen{$DTR\uparrow$}، بابنسکی \MDen{upward} و علائم درگیری مخچه‌ای.
\item برخلاف \MDen{PD} وضعیت خمیده ندارند.
\item در مراحل پیشرفته آتروفی میدبرین در \MDen{MRI} دیده می‌شود.
\item ۲۰\% بیماران به آمانتادین پاسخ می‌دهند.
\item تزریق بوتاکس برای رفع دیستونی گردن، بلفارواسپاسم.
\item طول عمر بیمار از شروع علائم ۷–۸ سال است.
\end{itemize}

% pic08.jpg = 20261012_224543235.jpg; handwritten page 8 (also a 9 at left).
\MDpage{۸}{pic08.jpg}{۱۴ ـ اختلالات حرکتی}
\noindent شمارهٔ حاشیهٔ چپ: ۹.
\section*{دیس‌کینزی‌ها}
\subsection*{انواع ترمور}
\begin{enumerate}[label=\alph*)]
\item ترمور استراحت (\MDen{Rest}): \MDen{PD} / سندرم‌های آکینتیک ـ ریجید.
\item ترمور \MDen{Action-Postural}:
\begin{itemize}
\item ترمور فیزیولوژیک.
\item ترمور فیزیولوژیک تشدیدشده: تیروتوکسیکوز / اختلالات اضطرابی / قطع الکل.
\item ترمور اسنشیال / سمپاتومیمتیک‌ها، لیتیوم، سدیم والپروآت.
\end{itemize}
\item ترمور \MDen{Intention}: ضایعات مخچه‌ای و \MDen{brainstem}: \MDen{MS} / بیماری عروقی / تومورها.
\end{enumerate}
\begin{itemize}
\item ترمور \MDen{Rest}: هنگامی که عضو درگیر در حالت رست باشد، تکیه کند یا زمانی که آویخته باشد.
\item ترمور \MDen{A-P}: هنگام حرکت ارادی یا حفظ عضو در برابر نیروی جاذبه.
\begin{itemize}
\item مثلاً کشیدن دست‌ها و نگهداری آن‌ها به صورت کشیده در جلوی تنه.
\end{itemize}
\item ترمور \MDen{Intention}: با نزدیک شدن عضو به هدف در یک حرکت هدفمند ایجاد می‌گردد: \MDen{Finger to nose}.
\end{itemize}
\section*{ترمور اسنشیال خوش‌خیم (\MDen{ET})}
\begin{itemize}
\item در زمان فعالیت ارادی، حفظ ارادی وضعیت عضو در مقابل نیروی جاذبه.
\item ۲ پیک سنی: ۱۵ سالگی و ۵۰ سالگی.
\item از \MDen{PD} بسیار شایع‌تر است.
\item الگوی وراثت \MDen{AD} در بیش از ۵۰\% بیماران.
\item وقتی فرد دست‌هایش را به جلوی تنه می‌گیرد تشدید می‌شود: کاراکتریستیک.
\item اعمال ظریف و صنایع‌دستی مختل.
\item دست‌خط به شدت مختل می‌شود.
\end{itemize}

% pic09.jpg = 20261012_224551392.jpg; handwritten page 9.
\MDpage{۹}{pic09.jpg}{۱۴ ـ اختلالات حرکتی}
\subsection*{ادامهٔ ترمور اسنشیال}
\begin{itemize}
\item فرکانس ۴–۸ هرتز.
\item شایع‌ترین محل: اندام‌های فوقانی / \MDen{$\tfrac{1}{4}$} موارد اندام تحتانی.
\item گردن، زبان، صدا، عضلات فک و صورت در موارد شدید.
\item وجود ترمور ایزوله در هنگام لبخند زدن به ضرر \MDen{ET}.
\item با بیماری‌های نورولوژیک دیگر همراهی ندارد.
\item احتمال وجود اختلالات خفیف شناختی، \MDen{gait} (\MDen{Unsteady gait}) یا آتاکسی راه رفتن وجود دارد.
\item گاهی حالت چرخ‌دنده‌ای (\MDen{Cogwheeling}) بدون ریجیدیتی دیده می‌شود.
\item میزان کم یا متوسط الکل به شدت ترمور را \MDen{$\downarrow$} می‌دهد.
\item در اغلب موارد سیر آهسته است.
\item درمان: بتابلاکر (پروپرانولول): در ۵۰\% بیماران پاسخ مناسب.
\item پریمیدون: گاهی کمک‌کننده است.
\item داروهای ضدپارکینسون اثری ندارند.
\item موارد شدید: تحریک عمقی یک‌طرفه یا دوطرفهٔ تالاموس.
\end{itemize}
\section*{ترمور مخچه‌ای}
\begin{itemize}
\item می‌تواند علامت مهم بیماری‌های مخچه.
\item ضایعات نیمکره‌های مخچه سبب ترمور \MDen{Intention} می‌شود.
\item \MDen{Finger to nose}: دیسمتری، هایپومتری، هایپرمتری.
\item \MDen{Heel to shin}.
\item یافته‌ها کلاسیک:
\begin{itemize}
\item ترمور \MDen{Intention}.
\item همراهی با دیسمتری.
\item عضلات محوری و پروگزیمال را درگیر می‌کند.
\end{itemize}
\item برای ترمور \MDen{Intention} درمان سود ندارد.
\end{itemize}

% pic10.jpg = 20261012_224558320.jpg; handwritten page 10.
\MDpage{۱۰}{pic10.jpg}{۱۴ ـ اختلالات حرکتی}
\begin{itemize}
\item ضایعات خط وسط و ورمیس مخچه: آتاکسی راه رفتن و \MDen{Titubation} سر و تنه.
\item ضایعات نیمکره‌های مخچه یا مسیرهای خروجی: آتاکسی اندام‌ها، ترمور \MDen{Intention}.
\item اتیولوژی ترمور \MDen{Intention}: ضایعات مخچه، ضایعات ساقهٔ مغز، \MDen{MS}، بیماری‌های عروقی و تومورها.
\end{itemize}
\section*{\MDen{Chorea}}
\begin{itemize}
\item حرکات سریع، مداوم، غیرارادی، نامنظم و بی‌هدف اندام‌ها با الگوی غیرقابل‌پیش‌بینی که از یک قسمت بدن به قسمت دیگر می‌رود.
\item بیشتر دیستال اندام‌ها را درگیر می‌کند.
\item دامنهٔ حرکات کم است.
\item علائم همراه: درهم‌کشیدن چهره / صداهای غیرطبیعی تنفسی / هایپوتونی.
\item بیمار معمولاً حرکات غیرطبیعی کره‌ای را در یک حرکت دلخواهٔ ارادی پنهان می‌کند.
\end{itemize}
\subsection*{کرهٔ سیدنهام}
\begin{itemize}
\item نادر / دنبال تب روماتیسمی، عفونت استرپتوکوک گروه \MDen{A} در کودکان.
\item با اختلالات رفتاری و روان‌پزشکی همراهی دارد.
\item گاهی آنتی‌بادی‌ها ضد گانگلیون‌های بازال همراهی دارد.
\item در دختر بچه‌ها شایع‌تر.
\item شروع تحت‌حاد دارد.
\begin{enumerate}[label=\alph*)]
\item هایپوتونی و ضعف عضلانی.
\item تغییرات رفتار همراه با علائم رفتاری و عاطفی.
\item یک‌طرفه یا دوطرفه.
\end{enumerate}
\item تشخیص بالینی است.
\item بعد از یک دورهٔ ۳–۶ ماهه به‌تدریج بهبود می‌یابد.
\item والپروآت مؤثر است: بهتر است از ابتدا شروع شود.
\item هالوپریدول در مرحلهٔ حاد مؤثر است.
\end{itemize}

% pic11.jpg = 20261012_224604013.jpg; handwritten page 11.
\MDpage{۱۱}{pic11.jpg}{۱۴ ـ اختلالات حرکتی}
\begin{itemize}
\item آنتی‌کولینرژیک همراه هالوپریدول ممنوع است.
\end{itemize}
\section*{همی‌بالیسم}
\begin{itemize}
\item حرکت پرتابی اندام‌های فوقانی و تحتانی در یک طرف.
\item ضایعات هستهٔ ساب‌تالامیک معمولاً طی استروک: همی‌بالیسم در سمت مقابل هستهٔ درگیر.
\item شبیه به کره است ولی بیشتر عضلات پروگزیمال اندام‌ها و کمربند شانه و لگن را درگیر می‌کند.
\item \MDen{RF}: سن بالا، \MDen{DM}، \MDen{HTN} \MDen{$\Leftarrow$} در واقع \MDen{RF}های استروک.
\item اگر ناشی از استروک باشد معمولاً طی چند ماه بهتر می‌شود.
\item داروهای نورولپتیک یا تترا‌بنازین کمک‌کننده هستند.
\end{itemize}
\section*{میوکلونوس}
\begin{itemize}
\item انقباض‌های کوتاه، ناگهانی، شوک‌مانند، منظم یا غیرمنظم در یک عضله یا گروهی از عضلات.
\item انواع:
\begin{itemize}
\item فیزیولوژیک: در شروع خواب.
\item پاتولوژیک:
\begin{enumerate}[label=\alph*)]
\item فوکال: فقط در یک ناحیهٔ آناتومیک.
\item سگمنتال: چند منطقهٔ آناتومیک مجاور هم.
\item مولتی‌فوکال: مناطق مختلف و زمان‌های گوناگون.
\item ژنرالیزه: از کورتکس مغز ناشی شده و در کل بدن مشاهده می‌شود.
\end{enumerate}
\end{itemize}
\item یکی از علائم کاراکتریستیک بیماری کروتزفلد ـ جاکوب می‌باشد.
\item درمان: سدیم والپروآت، کلونازپام، پریمیدون، پیراستام، لوتیراستام.
\end{itemize}

% pic12.jpg = 20261012_224620413.jpg; handwritten page 12.
\MDpage{۱۲}{pic12.jpg}{۱۴ ـ اختلالات حرکتی}
\section*{اسپاسم همی‌فاشیال}
\begin{itemize}
\item یک میوکلونوس فوکال شایع است که در یک نیمهٔ صورت (سمت ضایعه) رخ می‌دهد.
\item فشار به عصب فاشیال توسط عروق در محل خروج عصب از ساقهٔ مغز.
\item درمان: تزریق بوتاکس / کاربامازپین.
\end{itemize}
\section*{تیک}
\begin{itemize}
\item حرکات تکراری با الگوی مشخص (استروتایپ) در یک قسمت بدن.
\item بیمار می‌تواند به طور ارادی برای مدت کوتاهی تیک را مهار کند؛ برخلاف میوکلونوس.
\item مناطق شایع: صورت، سر، قسمت فوقانی تنه.
\item در کودکی شایع / اغلب گذرا / با افزایش سن برطرف می‌شود.
\end{itemize}
\subsection*{انواع}
\begin{enumerate}[label=\alph*)]
\item حرکتی ساده: پلک‌زدن مکرر، چین انداختن بینی، تکان‌دادن شانه همراه با کج‌کردن گردن.
\item حرکتی کمپلکس: چند گروه عضله درگیر هستند / دستکاری یا بوسیدن اشیا و بدن، نشان‌دادن حالت‌های مختلف چهره، خاراندن، لگدزدن.
\item دیستونیک: چرخاندن گردن، انقباض عضلات پلک شبیه بلفارواسپاسم؛ \MDunclear{تیک عرضی} آهسته، به صورت یک انقباض مداوم باشد.
\item صوتی ساده: صاف‌کردن گلو، سرفه، جیغ، ملچ‌ملوچ، خرناس کشیدن.
\item صوتی کمپلکس: بیان کلمات یا جملات خاص / عادی یا زشت یا بی‌معنی.
\end{enumerate}
\begin{itemize}
\item اندیکاسیون درمان تیک: اختلال در فعالیت‌های زندگی.
\item بیشتر تیک‌های سادهٔ خفیف نیاز به درمان ندارند.
\end{itemize}

% pic13.jpg = 20261012_224627317.jpg; handwritten page 13.
\MDpage{۱۳}{pic13.jpg}{۱۴ ـ اختلالات حرکتی}
\subsection*{ادامهٔ درمان تیک}
\begin{itemize}
\item داروها: آنتی‌دوپامینرژیک: پیموزاید، سولپرید، هالوپریدول، تترا‌بنازین.
\item درمان تیک‌های شدید: جراحی فانکشنال \MDen{DBS}.
\end{itemize}
\section*{سندرم ژیل دولاتوره (تورت)}
\begin{itemize}
\item تیک‌های حرکتی ـ کلامی متعدد و مادام‌العمر / وراثت \MDen{AD}.
\item قبل از ۱۸ سالگی شروع می‌شود.
\item با دوره‌های بهبودی و عود همراه است؛ در بزرگسالی شدت علائم کم می‌شود.
\item بیشتر فوقانی تنه، صورت، گردن و شانه‌ها.
\item الگوی تیک در طول زمان متغیر است.
\item ۱۵\% بیماران کلمات زشت بیان می‌کنند: \MDen{Coprolalia}.
\item بیماری‌های همراه: \MDen{ADHD}، اختلال وسواس جبری (\MDen{OCD}).
\begin{itemize}
\item \MDunclear{قابل تشدید است}: \MDen{SSRI}.
\end{itemize}
\end{itemize}
\section*{بیماری ویلسون}
\begin{itemize}
\item وراثت \MDen{AR} / اختلال در متابولیسم مس.
\item بیماری کبدی و اختلالات نورولوژیک.
\item جهش در ژن پروتئین انتقال‌دهندهٔ مس در کروموزوم ۱۳.
\begin{itemize}
\item مس از صفرا دفع نمی‌شود؛ در کبد و مغز تجمع می‌یابد؛ مرگ سلولی.
\end{itemize}
\item سن شروع ۵–۴۰ سالگی / بعد از ۴۰ سالگی نادر است.
\end{itemize}
\subsection*{علائم نورولوژیک}
\begin{enumerate}[label=\alph*)]
\item نوع دیستونیک:
\begin{itemize}
\item دیزآرتری، دیسفاژی دیستونیک، \MDunclear{راه رفتن دیستونیک}.
\item دیستونی عضلات بولبار و در نتیجه ریزش بزاق (\MDen{drooling}).
\item دیستونی عضلات صورت و زهرخند: \MDen{Risus Sardonicus}.
\item ترمور وضعیتی و فعالیتی در دیستال یا پروگزیمال.
\end{itemize}
\end{enumerate}

% pic14.jpg = 20261012_224632745.jpg; handwritten page 14.
\MDpage{۱۴}{pic14.jpg}{۱۴ ـ اختلالات حرکتی}
\subsection*{ادامهٔ علائم ویلسون}
\begin{itemize}
\item \MDen{Wined-beating tremor}: در نوع شدید و پروگزیمال، ترمور خاصی شبیه بال‌زدن ایجاد می‌کند.
\item به‌ندرت ترمور زبان در ابتدای بیماری تنها علامت است.
\end{itemize}
\begin{enumerate}[label=\alph*),start=2]
\item نوع آکینتیک ـ ریجید:
\begin{itemize}
\item درجات متفاوتی از برادی‌کینزی، ریجیدیتی و صورت ماسکه.
\item ترمور استراحت و ترمور وضعیتی شایع است.
\end{itemize}
\item نوع سودواسکلروتیک یا مخچه‌ای:
\begin{itemize}
\item شباهت زیادی با علائم بیماری \MDen{MS} دارد.
\item علائم مخچه‌ای: آتاکسی راه رفتن، دیزآرتری آتاکسیک، عدم هماهنگی حرکات.
\item ترمور \MDen{Intention} یا آتاکسیک شایع است؛ فرم شدید آن ترمور هولمز نامیده می‌شود.
\end{itemize}
\end{enumerate}
\subsection*{علائم غیرنورولوژیک}
\begin{itemize}
\item اختلالات رفتاری ـ شخصیتی، خلقی، شناختی، سایکوتیک.
\item اختلالات روان‌پزشکی، کبدی و درگیری چشمی.
\item ۲۰\% موارد با علائم روان‌پزشکی شروع می‌شود.
\item در نهایت همهٔ بیماران دچار علائم روان‌پزشکی می‌شوند.
% The word «علائمی» immediately before «علائم روان‌پزشکی» is crossed out.
\item شایع‌ترین عارضهٔ کبدی ویلسون: سیروز پست‌نکروتیک.
\end{itemize}
\subsection*{درگیری چشمی}
\begin{itemize}
\item حلقهٔ کایزر ـ فلشر (\MDen{KF}): شایع‌ترین یافتهٔ چشمی.
\begin{itemize}
\item رسوب مس در قرنیه.
\item همیشه دوطرفه است.
\item رنگ طلایی، قهوه‌ای، سبز.
\item روی دید تأثیری ندارد.
\item با اسلیت‌لمپ بهتر دیده می‌شود.
\end{itemize}
\end{itemize}

% pic15.jpg = 20261012_224641902.jpg; handwritten page 15.
\MDpage{۱۵}{pic15.jpg}{۱۴ ـ اختلالات حرکتی}
\begin{itemize}
\item کاتاراکت: به شکل گل آفتابگردان.
\begin{itemize}
\item دید را مختل نمی‌کند.
\item مشاهده با اسلیت‌لمپ.
\end{itemize}
\end{itemize}
\subsection*{سایر عوارض ویلسون}
\begin{itemize}
\item آنمی همولیتیک، اختلالات خونریزی‌دهنده، مشکلات کلیوی، استئومالاسی، ریکتز.
\end{itemize}
\section*{تشخیص بیماری ویلسون}
\begin{itemize}
\item سطح سرمی سرولوپلاسمین: در ۹۵\% بیماران \MDen{$\downarrow$} است.
\item مس ادرار ۲۴ ساعته: \MDen{$\uparrow$} می‌یابد.
\item معاینهٔ چشم: بررسی از نظر \MDen{KF}.
\item بیوپسی کبد: در موارد مشکوک، بررسی و اندازه‌گیری مس موجود در بافت کبد.
\item تصویربرداری با مس رادیواکتیو \MDen{Cu64}: مواردی که بیوپسی کبد امکان‌پذیر نیست.
\item \MDen{MRI} و \MDen{CT}:
\begin{itemize}
\item نمای صورت پاندا در \MDen{Midbrain}: \MDen{Face of giant panda}.
\item آتروفی کورتیکال.
\item تغییر سیگنال در مادهٔ سفید فرونتال و گانگلیون‌های بازال (به‌ویژه در پوتامن).
\item ضایعات تالاموس و ساقهٔ مغز.
\end{itemize}
\item اندازه‌گیری مس سرم ارزش تشخیصی ندارد.
\item حلقهٔ \MDen{KF} همواره همراه با علائم نورولوژیک دیده می‌شود.
\end{itemize}
\section*{درمان ویلسون}
\begin{itemize}
\item درمان دارویی مادام‌العمر.
\item محدودیت مس رژیم غذایی: شکلات، قارچ، صدف، جگر، دانه‌های روغنی.
\end{itemize}

% pic16.jpg = 20261012_224649701.jpg; handwritten page 16.
\MDpage{۱۶}{pic16.jpg}{۱۴ ـ اختلالات حرکتی}
\subsection*{ادامهٔ درمان ویلسون}
\begin{itemize}
\item جلوگیری از جذب مس در دستگاه گوارش.
\item ترکیب خوراکی سولفات یا استات روی: نباید همزمان با پنی‌سیلامین تجویز شود.
\item تترا‌تیومولیبدات: یک دورهٔ ۸ هفته‌ای در شروع درمان.
\item تخلیهٔ مس بدن.
\item پنی‌سیلامین: مؤثرترین داروی درمان ویلسون.
\begin{itemize}
\item داروی اصلی شروع درمان و درمان نگهدارنده.
\item \MDen{$\tfrac{2}{3}$} بیماران به آن پاسخ می‌دهند؛ سایر افراد ناتوانی شدید یا مرگ.
\item بهبود علائم طی چند هفته تا چند ماه.
\item در برخی بیماران به طور موقت علائم نورولوژیک تشدید می‌شوند.
\end{itemize}
\item تریانتین:
\begin{itemize}
\item تأثیر کمتر از پنی‌سیلامین.
\item عوارض آن کمتر از پنی‌سیلامین.
\item علائم نورولوژیک در شروع درمان تشدید می‌شوند.
\end{itemize}
\item پیوند کبد: در موارد شدید درگیری کبد؛ موجب بهبود در برخی علائم نورولوژیک می‌شود.
\end{itemize}
\section*{اختلالات حرکتی دارویی}
\subsection*{آکاتیزیا}
\begin{itemize}
\item بی‌قراری حرکتی سایکوموتور.
\item بیمار در حالت نشسته بی‌قرار است و دائماً دست و پایش را حرکت می‌دهد.
\item برای بهتر شدن شروع به راه رفتن می‌کند.
\item اتیولوژی: پارکینسونیسم دارویی و ایدیوپاتیک.
\item درمان: آنتی‌کولینرژیک‌ها، پروپرانولول، کلونازپام.
\end{itemize}

% pic17.jpg = 20261012_224656722.jpg; handwritten page 17.
\MDpage{۱۷}{pic17.jpg}{۱۴ ـ اختلالات حرکتی}
\section*{دیس‌کینزی تاردیو}
\begin{itemize}
\item حرکات غیرطبیعی و غیرارادی که پس از چند ماه تا چند سال درمان با نورولپتیک‌ها به وجود می‌آید.
\end{itemize}
\subsection*{علائم بالینی}
\begin{itemize}
\item شایع‌ترین علامت دیس‌کینزی کورئوآتتوز اوروفاشیال می‌باشد.
% «فاشیال» is corrected above the crossed-out ending of «اوروفاشیال».
\begin{itemize}
\item حرکات تکرارشوندهٔ لب، دهان، زبان.
\end{itemize}
\item در زنان مسن از میانسالی شایع‌تر است.
\item ۶۰\% بیماران با قطع داروی نورولپتیک، تدریجاً خوب می‌شوند.
\item در بیماران جوان دیس‌کینزی تاردیو به شکل سندرم دیستونیک تظاهر می‌یابد؛ اغلب با تورتیکولیس همراه با دیستونی سایر اندام‌ها.
\item فقط ۱۵\% بهبود می‌یابند.
\end{itemize}
\subsection*{درمان}
\begin{itemize}
\item اگر مصرف نورولپتیک به علت بیماری روان‌پزشکی است باید از داروهای نورولپتیک آتیپیک استفاده شود.
\item نوع اوروفاشیال: با تترا‌بنازین بهتر می‌شود؛ با آنتی‌کولینرژیک‌ها بدتر می‌شود.
\item دیستونی تاردیو: با آنتی‌کولینرژیک‌ها بهتر می‌شود.
\item درمان دیس‌کینزی تاردیو: تترا‌بنازین است.
\end{itemize}
\section*{سندرم نورولپتیک بدخیم (\MDen{NMS})}
\begin{itemize}
\item عارضهٔ داروهای بلاک‌کنندهٔ رسپتور دوپامین (\MDen{DRB}).
\item تب، ریجیدیتی شدید، ناپایداری اتونوم، اختلال هوشیاری.
\item \MDen{Lab}: لکوسیتوز، \MDen{$\uparrow$} کراتین‌کیناز سرم.
\item درمان: قطع دارو / تجویز آگونیست دوپامین / شل‌کنندهٔ عضلات (دانترولن).
\end{itemize}

\endgroup
\ifMotorStandalone\expandafter\enddocument\fi
